% SHARD-ID: AQM-78-QSA-FINITE-PHASE-WEYL-KINEMATICS
% SHARD-TITLE: Weyl-Heisenberg Kinematics From Finite Phase Spaces
% SHARD-SUMMARY: Consolidates the common Weyl-Heisenberg representation step used after a finite symplectic label space has been fixed.
% SHARD-SUMMARY: Separates finite-field, finite Frobenius-ring, sourced generalized-Heisenberg, and unresolved finite-ring cases.
% SHARD-SUMMARY: Records where uniqueness or projective-representation classification is locally sourced and where only Pauli/error-basis support is available.
% SHARD-KEYWORDS: quantum-system association, Weyl-Heisenberg, finite phase space, finite field, Frobenius ring, Stone-von Neumann, Clifford gap

\section{Weyl-Heisenberg Kinematics From Finite Phase Spaces}
\label{sec:qsa-finite-phase-weyl-kinematics}

\subsection{Status and Local Anchors}

This is QSA Step 14 from
\path{docs/quantum_system_associations/PLAN.md}.  It is a review/new
synthesis shard: it adds no source manifest, no convention entry, no run
artifact, and no new all-finite-ring theorem.  Its purpose is to isolate the
common representation-theoretic step used after earlier association workflows
have already produced a finite phase-label space.

The governing vocabulary is AQM-65 and \texttt{CONVENTIONS.md} convention
\texttt{(ap)}.  The finite-field arithmetic anchors are AQM-22, AQM-28,
AQM-60, AQM-76, and the finite symplectic field-label review AQM-71.  The
ring-status anchor is AQM-31.  The source manifests are
\path{references/symplectic_qecc/SOURCES.md} and
\path{references/heisenberg_weil/SOURCES.md}.  In particular, AQM-31 cites
Prasad, Beny--Crann--Lee--Park--Youn, Bekka, Lysenko,
Sidana--Kashyap, Gluesing-Luerssen--Pllaha, Solomon, and the finite
Clifford/Weil splitting sources at the line locators recorded there.

\subsection{The Project's Local Weyl Step}

The Weyl-Heisenberg step begins only after a finite label object and its
commutator pairing have already been fixed.  In the finite-field vector-space
case, the local display form is
\[
  E=Q\oplus P,\qquad \omega((q,p),(q',p'))=p(q')-p'(q),
\]
where \(Q\) is a finite-dimensional vector space over a finite field \(K\),
and \(P=Q^*\) or a chosen identified dual.  With
\[
  \psi_K(a)=
  \exp\!\left(\frac{2\pi i}{p}
    \operatorname{Tr}_{K/\mathbb F_p}(a)\right),
\]
the carrier is
\[
  \mathcal H=\ell^2(Q).
\]
For \(a\in Q\), define translations and phases by
\[
  T(q)|a\rangle=|a+q\rangle,\qquad
  R(p)|a\rangle=\psi_K(p(a))|a\rangle,
\]
and set
\[
  W(q,p)=T(q)R(p).
\]
Then
\[
  W(q,p)W(q',p')=\psi_K(p(q'))W(q+q',p+p')
\]
and
\[
  W(e)W(f)=\psi_K(\omega(e,f))W(f)W(e).
\]
This is the local finite-field construction used in AQM-22 for prime-field
qudit Paulis, AQM-28 for closed finite-residue sites, AQM-60 for
tangent-cotangent labels, and AQM-76 for the cotangent-first QSA workflow.
If an earlier workflow produces an abstract finite symplectic vector space
rather than an already polarized \(Q\oplus Q^*\), a choice of symplectic basis
or polarization is presentation data for writing this Schrodinger model.

This step is kinematical.  It supplies a projective representation, Weyl
operators, and the full finite matrix observable algebra on the carrier.
It does not choose a Hamiltonian, time evolution, de Rham CSS stabilizer,
state, Lagrangian, or code subspace.  AQM-22 T2, AQM-60, AQM-63, and AQM-76
all keep stabilizer or de Rham selections as later choices.

\subsection{Case 1: Prime-Field and Finite-Field Vector Spaces}

The prime-field case is fully local in the report.  AQM-22 uses
\[
  E_{\rm el}=\operatorname{Map}(\mathbb F_p,\mathbb F_p^2)
\]
and proves that the operators \(T_\xi R_\zeta\) generate the ambient
odd-\(p\) \(p\)-qudit Pauli group on the carrier
\(\ell^2(\mathbb F_p^{\mathbb F_p})\).  It also records the \(p=2\) phase
edge: the projective label space and commutator criterion remain the same,
but the conventional Hermitian Pauli group uses fourth roots of unity.

AQM-28 extends the same finite-field Weyl step to closed points of affine
finite-type \(\mathbb F_q\)-schemes.  A closed point \(x\) contributes
\[
  E_x=\kappa(x)^2,\qquad
  \mathcal H_x=\ell^2(\kappa(x)),
\]
and finite supports use direct sums of labels and tensor products of
carriers.  AQM-60 and AQM-76 use the same construction after replacing the
residue label \(K^2\) by the geometric tangent-cotangent label
\[
  T_{X/k,x}\oplus T^*_{X/k,x}.
\]
For \(m=\dim_{\kappa(x)}T_{X/k,x}\), the finite carrier dimension is
\(|\kappa(x)|^m\), as recorded in convention \texttt{(ao)} and reviewed in
AQM-76.

The sourced classification status is narrower than the construction.  The
finite-field Pauli/error-basis laws are anchored in
\path{references/symplectic_qecc/SOURCES.md}: Ashikhmin--Knill
\texttt{n9.tex:249-313}, Gottesman \texttt{Thesis.tex:1627-1681}, and
Bombin--Martin--Delgado \texttt{HomologicalErrorCorrection6.tex:1206-1259}.
For odd finite-dimensional Weyl systems, Gross is registered at
\texttt{poswig.tex:340-453} for Weyl operators and at
\texttt{poswig.tex:1966-2042} for a discrete Stone-von Neumann/Clifford proof.
Thus odd finite-field Clifford/Gaussian statements may be used only in that
sourced range.  The even case keeps the phase caveats already recorded in
AQM-22, AQM-31, and convention \texttt{(af)}.

\subsection{Case 2: Finite Frobenius-Ring Pauli and Stabilizer Cases}

AQM-31 records the finite Frobenius-ring status from
Gluesing-Luerssen--Pllaha and Sidana--Kashyap.  For a finite commutative
Frobenius ring \(R\) with generating character \(\chi\), the sourced Pauli
operators on \(\ell^2(R^n)\) are
\[
  X(a)v_x=v_{x+a},\qquad Z(b)v_x=\chi(b\cdot x)v_x,
\]
with
\[
  X(a)Z(b)X(a')Z(b')
  =
  \chi(b\cdot a')X(a+a')Z(b+b')
\]
and commutator controlled by
\[
  b\cdot a'-b'\cdot a.
\]
The local locators are
\path{references/heisenberg_weil/gluesing_luerssen_pllaha_1710_09884_source/StabCodesFrob5.tex},
lines 224--249, 286--323, 353--355, 412--423, and 469--504, and
\path{references/heisenberg_weil/sidana_kashyap_2202_00248_source/EAQECCs_over_rings_v15.tex},
lines 103--145 and 180--197.

The project already has one explicit Artin/Frobenius local proof in AQM-36.
For \(A=\mathbb F_3[t]/(\pi^e)\), it defines the top-coefficient character,
proves it is generating, constructs \(W_A(q,p)\) on \(\ell^2(A)\), proves the
projective product and commutator laws, and proves that the Weyl operators
form an operator basis of \(\operatorname{End}_{\mathbb C}(\ell^2(A))\).
This supports the nilpotent-sensitive thickened layer in conventions
\texttt{(ab)} and \texttt{(an)}.

The boundary is important.  These sources support Pauli/error bases,
character-symplectic commutation, and stabilizer self-orthogonality over the
stated finite Frobenius-ring scope.  They do not, in the current project,
provide a complete Clifford/Weil representation classification for every
finite commutative Frobenius ring, and they do not classify projective
representations over all finite rings.

\subsection{Case 3: Sourced Finite Generalized Heisenberg Cases}

AQM-31 also records Solomon's generalized finite Heisenberg source.  The
local support is
\path{references/heisenberg_weil/solomon_2501_00650_source/TotslipVersion5_Aug_25.tex},
lines 315--318, 904--981, 1127--1136, 1637--1649, and 1677--1718.  In that
setting, finite abelian groups \(A,B,C\) and a bilinear pairing
\[
  \lambda:A\times B\to C
\]
define a Heisenberg group with product
\[
  h(a,b,c)h(a',b',c')
  =
  h(a+a',b+b',c+c'+\lambda(a,b')).
\]
A character \(p:C\to\mathbb C^\times\) gives the Schrodinger representation
on functions on \(A\):
\[
  (h(a,b,c)\cdot f)(x)=p(\lambda(x,b)+c)f(x+a).
\]

This is a legitimate sourced finite Heisenberg kinematics whenever the
chosen finite groups, pairing, and character are exactly in that sourced
scope or are checked locally against it.  It is broader than finite fields in
shape, but it is not a license to treat every finite ring label module as
classified.  The QSA catalogue may use it as a finite generalized
Heisenberg slot only after recording the concrete \(A,B,C,\lambda,p\) data.

\subsection{Case 4: General Finite-Ring Questions and Gaps}

For a general finite commutative ring \(R\), the project does not currently
claim a complete Weyl-Heisenberg, Stone-von Neumann, Clifford, or Weil
classification.  The following remain gaps unless a later shard adds local
sources, conventions, or checked derivations.

\begin{itemize}
\item If \(R\) is not equipped with a certified generating character, the
translation/phase recipe \(R^n\oplus R^n\) does not yet have the
Frobenius-ring self-duality input used in AQM-31 and AQM-36.
\item If \(R\) is Frobenius, the Pauli/error-basis and stabilizer layer is
sourced as above, but full Clifford/Weil splitting or classification is
additional structure.  AQM-31 records finite Clifford splitting subtleties
from Hashimoto--Horibe--Hayashi, Korbelar--Tolar, and Galindo; it does not
turn them into a universal finite-ring classification.
\item Even-order and \(4\)-divisibility phenomena are genuine source-backed
warnings, not cosmetic phase choices.  AQM-31 records Lysenko's odd-exponent
scope and Galindo's finite-abelian splitting criterion from
\path{references/heisenberg_weil/SOURCES.md}.
\item Finite-ring database rows in convention \texttt{(an)} are
layer-labelled.  A residue row, a thickened-Frobenius row, and a blocked
quantisation row are different evidence states; none is a hidden
all-finite-ring theorem.
\end{itemize}

\subsection{Conclusion for the QSA Catalogue}

QSA Step 14 supplies a common final kinematical move:
\[
  \text{finite phase labels with commutator data}
  \quad\leadsto\quad
  \text{Weyl-Heisenberg projective representation and observable algebra}.
\]
For finite fields and the prime-field qudit cases, this move is already used
through AQM-22, AQM-28, AQM-60, AQM-71, and AQM-76.  For finite Frobenius
rings, it is available at the Pauli/stabilizer and checked Artin-Weyl level
described by AQM-31 and AQM-36.  For Solomon-style finite generalized
Heisenberg groups, it is available when the finite groups, pairing, and
central character are explicitly supplied.  General finite-ring
classification, full Clifford/Weil classification, and dynamics remain
questions, not conclusions.
